\begin{abstract}
Quantitative structure--activity relationship (QSAR) models for Ames mutagenicity are increasingly used in ICH M7 assessments,\cite{ICH_M7_2017} yet reported performance varies widely with evaluation methodology. We systematically evaluated 375 model configurations---seven algorithms, five feature representations, and three preprocessing scenarios---across five genotoxicity datasets: Ames ($n = 13{,}560$), in~vitro chromosome aberration ($n = 1{,}619$), in~vivo micronucleus ($n = 2{,}153$), and two negative-sampled variants. Using a single locked configuration (XGBoost, compact descriptors, raw SMILES) at three levels of evaluation rigor, Ames MCC fell from 0.670 (random-split CV) to 0.624 (scaffold-split CV) to 0.234 (drug$\to$industrial leave-domain-out). The split-method effect ($\Delta = 0.046$) was dwarfed by source heterogeneity ($\Delta = 0.390$): a naive source classifier achieved MCC = 0.608 without any structural information, reflecting a 19.5-fold prevalence gap between drug-sourced (58.3\%) and industrial-sourced (3.0\%) compounds. Threshold-tuning analysis attributed only part of this gap to prior shift, with oracle MCC plateauing at 0.27--0.36, well below in-domain performance. Algorithm robustness under domain shift was uneven: Random Forest maintained MCC = 0.325 in the hardest direction where XGBoost collapsed to 0.122. For in~vivo micronucleus, salt stripping enabled meaningful ranking (ROC-AUC = 0.860, PR-AUC = 0.237) but classification remained unreliable (MCC 95\% CI crossing zero). We recommend multi-level evaluation with source-awareness as a reporting standard for Ames mutagenicity QSAR.
\end{abstract}
